% Ingeniería de Requisitos (ISR-401) · UTEQ
% TA-IND-01 · Caso hipotético de cobro predial
% Borrador de trabajo: alcance, actores, RF y requisitos de calidad RNF.
% Las decisiones técnicas y la redacción requieren revisión del estudiante.
\documentclass[10pt,onecolumn]{IEEEtran}

\usepackage[T1]{fontenc}
\usepackage[utf8]{inputenc}
\usepackage[spanish,es-tabla]{babel}
\usepackage{lmodern}
\usepackage{microtype}
\usepackage{csquotes}
\usepackage{booktabs}
\usepackage{tabularx}
\usepackage{array}
\usepackage{enumitem}
\usepackage[hidelinks]{hyperref}
\usepackage[backend=biber,style=ieee,sorting=none]{biblatex}
\addbibresource{referencias.bib}
\DefineBibliographyStrings{spanish}{bibliography={Referencias}}
\urlstyle{same}

\setlength{\parindent}{0.65em}
\setlength{\parskip}{2pt}
\setlength{\emergencystretch}{2em}
\renewcommand{\arraystretch}{1.22}
\newcolumntype{Y}{>{\raggedright\arraybackslash}X}
\newcommand{\rf}[2]{\par\smallskip\noindent\textbf{#1.} #2\par}
\newcommand{\prueba}[1]{\noindent{\footnotesize\textit{Comprobación propuesta:} #1}\par}

\title{Requisitos iniciales para un sistema de cobro de impuestos prediales en línea\\
\large TA-IND-01 · Avance técnico 2 (borrador)}
\author{Cedeño Avila Winston Damian\\
{\small Universidad Técnica Estatal de Quevedo · Facultad de Ciencias de la Computación}\\
{\small Ingeniería de Software · Cuarto nivel «B» · Ingeniería de Requisitos (ISR-401)}\\
{\small 2 de noviembre de 2026}}

\begin{document}
\maketitle

\begin{abstract}
Este segundo avance analiza un caso de clase: un GAD cantonal necesita un portal para consultar y pagar impuestos prediales en línea. Se mantienen el alcance preliminar, seis partes interesadas y cinco requisitos funcionales del primer borrador. Ahora se incorporan cinco requisitos no funcionales con objetivos numéricos, condiciones de medición y métodos de verificación propuestos, utilizando ISO/IEC/IEEE 29148:2018 e ISO/IEC 25010:2023 como referencias \cite{iso29148,iso25010}. El GAD es hipotético: todavía no hay entrevistas ni resultados de pruebas reales.
\end{abstract}

\section{Situación y alcance inicial}
El problema no consiste solamente en mostrar el valor de una deuda. Cuando alguien realiza un pago, el municipio necesita recibir una confirmación confiable antes de registrar la obligación como cancelada. Si esa confirmación falla o llega dos veces, el portal no debería mostrar un resultado incorrecto.

Para este ejercicio se propone que el portal permita consultar obligaciones de un predio, preparar el valor de pago, comunicarse con una pasarela local, confirmar el resultado y entregar un comprobante. También se considera la integración solicitada con el Servicio de Rentas Internas (SRI), aunque \textbf{no se presupone que exista una API disponible}: antes de implementarla se necesitaría revisar sus condiciones de acceso. Los valores de deuda, el catastro y las reglas municipales seguirían a cargo del GAD.

\textbf{Límite de este avance:} se han especificado funciones y objetivos de calidad para un caso hipotético. Los valores medibles son \textbf{metas académicas propuestas}, no cifras obtenidas de un GAD. No se ha desarrollado el sistema, ejecutado pruebas ni confirmado convenios con proveedores. Las restricciones y dependencias se revisarán después.

\section{Partes interesadas y necesidades}
Una parte interesada (\emph{stakeholder}) es una persona, grupo o entidad que usa, administra o afecta el sistema. Las necesidades de la Tabla~\ref{tab:actores} son hipótesis iniciales derivadas del caso; deberán confirmarse si se trabaja con un GAD real.

\begin{table}[htbp]
\centering
\small
\caption{Actores y necesidades preliminares del caso}
\label{tab:actores}
\begin{tabularx}{\linewidth}{@{}p{3.85cm}Y@{}}
\toprule
\textbf{Actor} & \textbf{Necesidad preliminar}\\
\midrule
Contribuyente & Consultar las obligaciones de su predio, pagar en línea y obtener una constancia.\\
Tesorería del GAD & Revisar pagos confirmados y evitar registros duplicados o sin respaldo.\\
Catastro y Rentas & Proporcionar claves catastrales, saldos y reglas de cálculo aprobadas por el municipio.\\
Pasarela de pagos local & Recibir una orden con identificador y monto y devolver el resultado de la transacción.\\
SRI & Establecer, cuando corresponda, las condiciones autorizadas de consulta o intercambio de datos tributarios.\\
Unidad de tecnología del GAD & Administrar permisos, integraciones y errores del portal.\\
\bottomrule
\end{tabularx}
\end{table}

\section{Requisitos funcionales propuestos}
Un \textbf{requisito funcional (RF)} describe una acción o respuesta observable del sistema. Para evitar frases vagas, cada enunciado utiliza «El sistema deberá» e indica la situación en la que ocurre la función. Los ejemplos de comprobación son \textbf{pruebas previstas}, no pruebas ya realizadas \cite{iso29148}.

\rf{RF-01 · Consulta de obligaciones}{El sistema deberá mostrar al contribuyente autenticado las obligaciones pendientes del predio asociado a su cuenta cuando seleccione la clave catastral y el período fiscal, incluyendo concepto, valor y estado de cada obligación.}
\prueba{Con un predio de prueba asociado al usuario y otro ajeno, comprobar que únicamente aparezcan los datos autorizados y que los importes coincidan con los registros de prueba del GAD.}

\rf{RF-02 · Cálculo de monto}{El sistema deberá calcular el total exigible de las obligaciones que el contribuyente seleccione, aplicando los valores, descuentos o recargos vigentes que entregue el GAD para la fecha de consulta.}
\prueba{Usar casos de prueba con montos y reglas previamente definidos por el municipio; comparar el total calculado con el resultado esperado.}

\rf{RF-03 · Pago y confirmación}{El sistema deberá crear una orden de pago con identificador único y enviarla a la pasarela local cuando el contribuyente confirme la operación. Solo marcará la orden como pagada al recibir una respuesta válida que coincida con su identificador y monto; una confirmación repetida no deberá registrar otro pago.}
\prueba{Simular una respuesta correcta, otra repetida y otra con monto distinto; comprobar que la orden se cancele una sola vez y que las respuestas inválidas no se acepten.}

\rf{RF-04 · Comprobante de pago}{El sistema deberá permitir descargar un comprobante PDF cuando la orden haya quedado confirmada como pagada; el comprobante incluirá predio, valor, fecha y referencia de transacción.}
\prueba{Comprobar que una orden confirmada genere el documento y que una orden pendiente o rechazada no genere constancia de pago realizado.}

\rf{RF-05 · Integración condicionada con el SRI}{El sistema deberá consultar la información tributaria necesaria mediante un mecanismo autorizado del SRI, siempre que exista un servicio habilitado y un convenio aplicable al GAD. Si el servicio no está disponible, informará esa condición sin modificar la deuda municipal.}
\prueba{Cuando se defina la interfaz real, probar una respuesta válida y una indisponibilidad controlada en un entorno autorizado; verificar que ninguna altere las obligaciones prediales.}

\section{Requisitos no funcionales y calidad}
A diferencia de los RF, que indican \emph{qué} operaciones realiza el portal, los \textbf{requisitos no funcionales (RNF)} señalan \emph{con qué nivel de calidad} debe funcionar. ISO/IEC 25010:2023 establece nueve características de calidad: adecuación funcional, eficiencia de desempeño, compatibilidad, capacidad de interacción, fiabilidad, seguridad, mantenibilidad, flexibilidad y seguridad operacional (\emph{safety}) \cite{iso25010}. Para este caso se seleccionan cinco características diferentes por su relación con consultas, pagos y atención ciudadana.

\rf{RNF-01 · Eficiencia de desempeño}{Con 100 usuarios virtuales concurrentes durante 30 minutos y una base de ensayo de 25\,000 predios \emph{simulados}, el 95\,\% de las consultas de obligaciones deberá completarse en 3 segundos o menos.}
\prueba{Métrica: percentil 95 (p95) del tiempo de respuesta, en segundos; objetivo: p95 $\leq 3$ s. Método: prueba de carga propuesta con JMeter, registrando únicamente consultas válidas y tiempos de respuesta del servidor.}

\rf{RNF-02 · Fiabilidad}{Durante un piloto de 30 días, el portal deberá alcanzar al menos un 99,5\,\% de consultas sintéticas exitosas al servicio de obligaciones prediales, con una solicitud automatizada cada minuto desde un punto externo.}
\prueba{Métrica: (consultas sintéticas correctas / consultas programadas)$\times100$; objetivo: $\geq 99{,}5$\,\%. Método: monitorización propuesta por 30 días; una consulta será correcta solo si devuelve el estado HTTP y los datos esperados.}

\rf{RNF-03 · Seguridad}{El portal deberá rechazar el 100\,\% de un conjunto de 150 intentos de consulta de predios o comprobantes sin autorización, usando cuentas y predios de ensayo con permisos conocidos.}
\prueba{Métrica: (intentos indebidos rechazados / 150)$\times100$; objetivo: 100\,\% en la batería definida. Método: pruebas automatizadas de autorización, incluyendo acceso directo a identificadores ajenos. Aprobar estos casos no equivale a demostrar ausencia total de vulnerabilidades.}

\rf{RNF-04 · Capacidad de interacción}{Al menos 18 de 20 ciudadanos que utilicen por primera vez un prototipo funcional deberán poder consultar el saldo de un predio de ensayo en 120 segundos o menos, sin recibir ayuda, desde un teléfono móvil con conexión estable.}
\prueba{Métrica: participantes que terminan dentro de 120 s / 20 $\times100$; objetivo: $\geq 90$\,\%. Método: prueba de usabilidad propuesta con cronómetro, instrucciones iguales para todos y registro anónimo del éxito o error de cada participante.}

\rf{RNF-05 · Compatibilidad / interoperabilidad}{Cuando exista un contrato de integración aprobado y un ambiente de pruebas de la pasarela local, el portal deberá interpretar correctamente las 50 respuestas válidas del conjunto de homologación, respetando sus identificadores, importes y estados.}
\prueba{Métrica: mensajes válidos procesados conforme al contrato / 50 $\times100$; objetivo: 100\,\% del conjunto definido. Método: pruebas de integración propuestas en entorno de homologación con mensajes preparados según la documentación real del proveedor; la interfaz todavía no se conoce.}

\subsection{Justificación de las metas}
El tiempo máximo de 3 s pretende que una consulta no obligue al ciudadano a esperar demasiado; los 100 usuarios y 25\,000 predios son valores de prueba para este ejercicio, no datos oficiales ni una estimación validada de la carga real. La disponibilidad del 99,5\,\% se plantea como punto de negociación. Para autorizaciones e interoperabilidad se exige un 100\,\% \textbf{solo en los casos de ensayo definidos}. La meta de usabilidad equivale a que al menos 18 de 20 personas completen la consulta sin asistencia. Estos valores deberán ajustarse con los responsables del GAD cuando se conozcan los recursos y las necesidades reales.

\section{Aspectos pendientes de revisión}
Aún deben formularse y justificarse las tres restricciones y las dos dependencias pedidas por la guía. Además, antes de convertir esta propuesta en una especificación aprobada, sería necesario confirmar las reglas municipales de cobro, el contrato de la pasarela, la posibilidad real de integración con el SRI y los umbrales de calidad con los responsables del GAD. \textbf{No se han aplicado pruebas de carga, seguridad, usabilidad ni disponibilidad}: por ahora solamente están definidos sus métodos de verificación.

\section*{Referencias}
\printbibliography[heading=none]

\noindent\footnotesize\textbf{Nota sobre asistencia de IA.} ChatGPT (OpenAI, GPT-6) se utilizó para elaborar el borrador técnico, proponer métricas de calidad y dar soporte a la edición LaTeX. Las metas numéricas no son mediciones reales y deben revisarse por el estudiante; no se realizaron entrevistas ni pruebas del sistema.

\end{document}
